\documentclass[10pt,a4paper]{amsart}
\usepackage[margin=19mm]{geometry}
\usepackage{amsmath,amssymb,mathtools}
\usepackage[scr=rsfs,cal=euler]{mathalfa}
\usepackage{xfrac}
\usepackage[unicode,hidelinks,bookmarks=false]{hyperref}
\newcommand{\ie}{i.e.\ }
\newcommand{\cf}{cf.\ }
\newcommand{\resp}{respectively\ }
\newcommand{\Subsection}{\subsection}
\providecommand{\nobreakdash}{\nobreak\hbox{-}}
\setlength{\emergencystretch}{2em}
\multlinegap0pt
\usepackage{bm}
\usepackage{bbm}
\usepackage{dsfont}
\usepackage{xspace}
\usepackage{cancel}
\usepackage{mathtools}
\usepackage{amsthm}
\makeatletter
\@ifundefined{theorem}{\newtheorem{theorem}{Theorem}}{}
\@ifundefined{lemma}{\newtheorem{lemma}[theorem]{Lemma}}{}
\@ifundefined{remark}{\newtheorem{remark}[theorem]{Remark}}{}
\makeatother
\def\R{\mathbb R}
\def\msD{\mathscr D}
\newcommand{\Eb}{\mathbf{E}}
\newcommand{\Ker}{\mathrm{Ker}}
\newcommand{\Tb}{\mathbf{T}}
\newcommand{\Tr}{\mathrm{Tr}}
\newcommand{\eb}{\mathbf{E}}
\newcommand{\ext}{\mathrm{ext}}
\newcommand{\mfn}{\mathfrak{n}}
\newcommand{\msE}{\mathscr{E}}
\newcommand{\rd}{\mathrm{d}\!\,}
\title[Quantum $SL^+(N,\mathbb{R})$ as a locally compact quantum gr]{Quantum $SL^+(N,\mathbb{R})$ as a locally compact quantum group}
\author{K. De Commer, G. Schrader, A. Shapiro, C. Voigt}
\date{Source: arXiv:2512.21579v1 (CC BY 4.0)}
\begin{document}
\maketitle

\noindent\emph{Source and licence.} Quantum $SL^+(N,\mathbb{R})$ as a locally compact quantum group. Version arXiv:2512.21579v1 (CC BY 4.0), distributed under the Creative Commons Attribution 4.0 International licence, as recorded in the arXiv licence metadata for this exact version and shown on the abstract page https://arxiv.org/abs/2512.21579v1. Full rights notice: licenses/arxiv-2512.21579-CC-BY-4.0.txt.\par\smallskip

\noindent\textit{Editorial scope.} Counted unit: the Lemma whose source label is \texttt{LemInvWeightQT} in the article (source0.tex lines 4483--4492, source label \texttt{LemInvWeightQT}), delivered together with the article's own complete original proof of it (source0.tex lines 4493--4555, one \texttt{proof} environment of 63 lines). No mathematics is written, repaired or re-derived: statement and proof are the article's own text line for line. Required context retained uncounted: EqTransDual (lines 4266--4268); RemWeight (lines 4463--4469); EqEqualFormPos (lines 6649--6653). Every definition, display and label of those lines is retained unchanged. Omitted from this package and not redistributed: 1--4265, 4269--4462, 4470--4482, 4556--6648, 6654--7246. Nothing inside a retained excerpt is omitted or altered: every retained body is delivered exactly as the article's lines are written, so no omission receipt of any kind accompanies this package. Citations: the retained excerpts carry no citation command at all, so no bibliography entry of the article is transcribed and no reference is redistributed. Reference adaptation: none was needed and none was made; every reference inside the delivered unit resolves inside it, so no unresolved reference or citation is produced. Printed slips kept literal: no printed word, symbol, hypothesis or number of the counted statement or of its proof was altered. Layout and class: the article's own class is \texttt{amsart} with packages including xcolor, amscd, amsmath, amssymb, amsthm, yfonts, mathrsfs, hhline, tikz, geometry, xy, parskip, tikz-cd, multirow; the wrapper is the lane's pinned \texttt{amsart} editorial wrapper at 10pt on A4, which restates the theorem-like environments the retained text needs and adds the article's own macro definitions that the retained text needs (\texttt{\textbackslash Eb}, \texttt{\textbackslash Ker}, \texttt{\textbackslash R}, \texttt{\textbackslash Tb}, \texttt{\textbackslash Tr}, \texttt{\textbackslash eb}, \texttt{\textbackslash ext}, \texttt{\textbackslash mfn}, \texttt{\textbackslash msD}, \texttt{\textbackslash msE}, \texttt{\textbackslash rd}), with extra wrapper packages (bm, bbm, dsfont, xspace, cancel, mathtools). The delivered numbering is the unit's own. Assets: no retained span contains a figure environment, a graphics-inclusion command, a PDF-page inclusion, a table or a TikZ picture; no image file is copied into this package, the package declares no asset dependency and no image file is redistributed with it. Licence: version arXiv:2512.21579v1 of this article is distributed under the Creative Commons Attribution 4.0 International licence, as recorded in the arXiv licence metadata for this exact version and shown on the abstract page https://arxiv.org/abs/2512.21579v1. A full rights notice is delivered at licenses/arxiv-2512.21579-CC-BY-4.0.txt. Characters: every retained excerpt is pure ASCII, so the delivered engine, XeLaTeX with its default Latin Modern text font, reports no missing character.
\begin{equation}\label{EqTransDual}
(\theta_{\hat{v}}f)(\hat{w}) = f(\hat{w}-\hat{v}),\qquad f\in L^2(\hat{V}),\hat{v},\hat{w}\in \hat{V}.
\end{equation}

\begin{remark}\label{RemWeight}
If $W \subseteq V$ is any vector space complementary to $Z$, so $V = Z\oplus W$, we easily see that 
\[
\msE_{\hat{W}}^{\hbar} (x) = \left(\int_{\hat{W}}^{\hbar}x\right)1,\qquad x\in L_{\hbar}(W)_+ \subseteq L_{\hbar}(V)_+,
\]
where we consider $\hat{W}$ as the dual space of $W$ upon restriction.
\end{remark} 

\begin{equation}\label{EqEqualFormPos}
\omega_{\xi,\xi}(x) = \begin{cases} \|\mathbf{x}^{1/2}\xi\|^2&\text{ if } \xi \in \msD(\mathbf{x}^{1/2}),\\
\infty &\text{ if } \xi \notin\msD(\mathbf{x}^{1/2}).
\end{cases}
\end{equation}

\begin{lemma}\label{LemInvWeightQT}
Assume $V$ is symplectic. Let $Z \subseteq V$, put $\hat{W} = \Ker(Z) \subseteq \hat{V}$, and choose a subspace $W \subseteq V$ such that
\[
V = Z\oplus W
\]
as vector spaces. Assume $z\in Z,w\in W$. Then there exists $a \in L_{\hbar}(W) \subseteq L_{\hbar}(V)$ such that  
\begin{equation}\label{EqFact}
\int_{\hat{Z}}^{\hbar,\eb(z)} =  \left(\Tr_{\Eb(z+w)}(a-a^*)\right)_{\mid L_{\hbar}(Z)_+}.
\end{equation}
\end{lemma} 

\begin{proof}
In the following, we again identify $\hat{W}$ with the dual of $W$ by the restriction map.

Consider the $\sigma$-weakly continuous one-parameter group of automorphisms 
\[
\gamma_t\colon  L_{\hbar}(W) \rightarrow L_{\hbar}(W),\qquad x \mapsto \Eb(z)^{it}x \Eb(z)^{-it}.
\]
This is clearly well-defined, and scales the weight $\int_{\hat{W}}^{\hbar,\Eb(w)}$:
\[
\int_{\hat{W}}^{\hbar,\Eb(w)} \circ \gamma_t = e^{2\pi\hbar t(w,z)}\int_{\hat{W}}^{\hbar,\Eb(w)}.
\]
Taking a non-zero positive $c \in \mfn_{\int_{\hat{W}}^{\hbar}}$ and putting 
\[
b = \int_{\R} e^{-t^2/2} \gamma_t (\widetilde{c})\rd t,\qquad \widetilde{c} = \chi_{[0,1]}(\Eb(w))c \chi_{[0,1]}(\Eb(w)),
\]
we see that we can find a non-zero $b \in L_{\hbar}(W)$ which is square $\int_{\hat{W}}^{\hbar,\Eb(w)}$-integrable, with $\Eb(w)^{1/2}b$ bounded, and with $b$ and  $\Eb(w)^{1/2}b$ analytic for $\gamma_t$. Normalize $b$ such that 
\begin{equation}\label{EqNormb}
\int_{\hat{W}}^{\hbar,\Eb(w)}bb^* =1,
\end{equation}
and put then 
\[
a = \gamma_{i/2}(b) \in L_{\hbar}(W).
\]
Then also $\Eb(w)^{1/2}a$ is bounded, and $\Eb(w)^{1/2}a$ is analytic for $\gamma_t$.

Now for $y \in L_{\hbar}(V)$ and $\Tb \in L_{\hbar}(V)_+^{\ext}$, let us write $y^*\cdot \Tb \cdot y \in L_{\hbar}(V)_+^{\ext}$ for the element 
\[
\omega(y^*\cdot \Tb\cdot y) = \omega(y^*-y)(\Tb),\qquad \omega \in M_*^+.
\]
We then claim that 
\begin{equation}\label{EqEqualPosOp}
\msE_{\hat{W}}^{\hbar}(a^*\cdot \eb(z+w)\cdot a) = \eb(z). 
\end{equation}
Indeed, take $\xi \in \msD(\eb(z)^{1/2}) \subseteq L^2(\hat{V})$. Then we compute, using the notation \eqref{EqTransDual}, that
\begin{eqnarray}
\omega_{\xi,\xi}(\eb(z)) 
&=& \|\eb(z)^{1/2}\xi\|^2 \\
&\underset{\eqref{EqNormb}}{=}&
\left(\int_{\hat{W}}^{\hbar,\Eb(w)} bb^*\right) \|\eb(z)^{1/2}\xi\|^2\\
&=& \left(\int_{\hat{W}}^{\hbar} (\Eb(w)^{1/2}b)^*(\Eb(w)^{1/2}b)\right)  \|\eb(z)^{1/2}\xi\|^2\\
&\underset{\textrm{Remark } \ref{RemWeight}}{=}& 
\omega_{\Eb(z)^{1/2}\xi,\Eb(z)^{1/2}\xi}\left(\msE_{\hat{W}}^{\hbar} \left((\Eb(w)^{1/2}b)^*(\Eb(w)^{1/2}b)\right) \right)\\
&=& \int_{\hat{W}} \|\Eb(w)^{1/2}b\theta_{\hat{w}}^* \eb(z)^{1/2}\xi\|^2\rd \hat{w}\\
&=& \int_{\hat{W}} \|\Eb(w)^{1/2}b \eb(z)^{1/2}\theta_{\hat{w}}^*\xi\|^2\rd \hat{w}\\
 &=& \int_{\hat{W}} \|\Eb(w)^{1/2} \Eb(z)^{1/2}a \theta_{\hat{w}}^*\xi\|^2 \rd\hat{w} \\
 &=& \int_{\hat{W}} \| \Eb(z+w)^{1/2}a \theta_{\hat{w}}^*\xi\|^2 \rd\hat{w} \\
&=& 
\omega_{\xi,\xi}\left(\int_{\hat{W}}^{\hbar} a^*\cdot \eb(z+w)\cdot a\right).
\end{eqnarray}
Writing the left hand side of \eqref{EqEqualPosOp} as $\Tb$, this shows by \eqref{EqEqualFormPos} that $\Eb(z)^{1/2} \subseteq \Tb^{1/2}$, which by self-adjointness of both sides is sufficient to conclude that they are equal. Hence equality holds in \eqref{EqEqualPosOp}.




Choose now $x\in L_{\hbar}(Z)$. Then we compute that 
\begin{eqnarray*}
\Tr_{\Eb(z+w)}(axx^*a^*) &=& \int_{\hat{V}}^{\hbar} x^*a^* \cdot  \eb(z+w)\cdot ax\\
&=& \int_{\hat{Z}}^{\hbar} \msE_{\hat{W}}^{\hbar} \left(x^*\cdot (a^* \cdot \eb(z+w)\cdot a)\cdot x\right)\\
&=& \int_{\hat{Z}}^{\hbar}x^*\cdot \left( \msE_{\hat{W}}^{\hbar} \left( a^*\cdot \eb(z+w)\cdot a\right)\right)\cdot x\\
&\underset{\eqref{EqEqualPosOp}}{=}& \int_{\hat{Z}}^{\hbar}x^*\cdot \eb(z)\cdot x\\
&=& \int_{\hat{Z}}^{\hbar,\Eb(z)}xx^*.
\end{eqnarray*} 
\end{proof}

\end{document}
